\section{Notations}\label{sec:prelim}
We first recall three basic operations. 
Given a permutation $\pi \in \mathfrak{S}_d$,  an integer $1 \le p \le n$ and tensors $T = u_1 \otimes \cdots \otimes u_d \in \mathbb{K}^{n_1} \otimes_{\mathbb{K}} \cdots \otimes_{\mathbb{K}} \mathbb{K}^{n_d}$,  $S = v_1 \otimes \cdots \otimes v_d \in \mathbb{K}^{m_1} \otimes_{\mathbb{K}} \cdots \otimes_{\mathbb{K}} \mathbb{K}^{m_d}$,  $F = f_{1} \otimes \cdots \otimes f_p  \in (\mathbb{K}^{n_{\pi(1)}})^\ast \otimes_{\mathbb{K}} \cdots \otimes_{\mathbb{K}} (\mathbb{K}^{n_{\pi(p)}})^\ast$,  the \emph{Kronecker product} of $T$ and $S$ is 
\[
T \boxtimes_{\mathbb{K}} S = (u_1\otimes v_1) \otimes \cdots \otimes (u_d \otimes v_d) \in \mathbb{K}^{n_1 m_1} \otimes_{\mathbb{K}} \cdots \otimes_{\mathbb{K}} \mathbb{K}^{n_d m_d}.
\]
We denote by $\langle T,  F \rangle $ or $\langle F,  T \rangle$ the tensor in $\mathbb{K}^{n_{\pi(p+1)}} \otimes_{\mathbb{K}} \cdots \otimes_{\mathbb{K}} \mathbb{K}^{\pi(n)}$ obtained by \emph{contracting $T$ with $F$}:
\[
\langle T,  S \rangle \coloneqq
\left( \prod_{i=1}^p f_i (u_{\pi(i)}) \right) u_{\pi(p+1)} \otimes \cdots \otimes u_{\pi(n)}. 
\]
The \emph{direct sum} of $T$ and $S$ is 
\[ 
T \oplus S \coloneqq (u_1,0)\otimes \cdots \otimes (u_d,0)+(0,v_1)\otimes \cdots \otimes (0,v_d) \in \mathbb{K}^{n_1 + m_1} \otimes_{\mathbb{K}} \cdots \otimes_{\mathbb{K}} \mathbb{K}^{n_d + m_d},
\]
where $(u,0)$ (resp.  $(0,v)$) denotes the image of $u \in \mathbb{K}^{n}$ (resp.  $v \in \mathbb{K}^{m}$) in $\mathbb{K}^{n + m}$ via the natural embedding. We extend these operations for general tensors by linearity.  
\subsection{Ranks of tensors}
For $S \in \mathbb{K}^{n_1} \otimes_{\mathbb{K}} \cdots \otimes_{\mathbb{K}} \mathbb{K}^{n_d}$,  if there exists some partition $\pi \in \mathfrak{S}_d$,  integer $1 \le p \le d-1$,  tensors $S_1 \in \mathbb{K}^{n_{\pi(1)}} \otimes_{\mathbb{K}} \cdots \otimes_{\mathbb{K}}  \mathbb{K}^{n_{\pi(p)}}$ and $S_2 \in \mathbb{K}^{n_{\pi(p+1)}} \otimes_{\mathbb{K}} \cdots \otimes_{\mathbb{K}} \mathbb{K}^{n_{\pi(d)}}$ such that
\[
S  = S_1 \otimes S_2,
\]
then we say $S$ has \emph{partition rank one}.  In particular,  if $p = 1$,  we say $S$ has \emph{slice rank one}.  Here the equality is understood via the natural isomorphism 
\[
\mathbb{K}^{n_1} \otimes_{\mathbb{K}} \cdots \otimes_{\mathbb{K}} \mathbb{K}^{n_d} \simeq 
\mathbb{K}^{n_{\pi(1)}} \otimes_{\mathbb{K}} \cdots \otimes_{\mathbb{K}}  \mathbb{K}^{n_{\pi(p)}}  \otimes_{\mathbb{K}}   \mathbb{K}^{n_{\pi(p+1)}} \otimes_{\mathbb{K}} \cdots \otimes_{\mathbb{K}}  \mathbb{K}^{n_{\pi(d)}}. 
\]
Given $T \in \mathbb{K}^{n_1} \otimes_{\mathbb{K}}  \cdots \otimes_{\mathbb{K}} \mathbb{K}^{n_d}$,  we denote 
\begin{align}
\pr_{\mathbb{K}}(T) &\coloneqq \min \left\lbrace
r \in \mathbb{N}: T = \sum_{j=1}^r T_j,\; T_j \text{~has partition rank one}
\right\rbrace,  \label{eq:PR} \\
\sr_{\mathbb{K}}(T) &\coloneqq \min \left\lbrace
r \in \mathbb{N}: T = \sum_{j=1}^r T_j,\; T_j \text{~has slice rank one}
\right\rbrace.  \label{eq:SR}
\end{align}
